\documentclass[11pt,a4paper]{article}
% =====================================================================
%  Shared preamble for the Part V lecture notes (lectures 19-22).
%
%  These four lectures sit outside Geron, so the notes ARE the primary
%  source and are examined on the same terms as the textbook parts.
%  Everything here is chosen to make that readable on paper: one column,
%  generous measure, and the same six-colour palette as the slides so a
%  student moving between the two is not re-learning what red means.
%
%  Build with pdflatex (twice, for the toc and the refs):
%      cd notes && latexmk -pdf lecture-19.tex
%  or  make -C notes
% =====================================================================
\usepackage[T1]{fontenc}
\usepackage[utf8]{inputenc}
\usepackage{newtxtext,newtxmath}      % Times-like: prints well, reads well
\usepackage[scaled=0.86]{inconsolata} % code, at a size that sits beside Times
\usepackage{microtype}
% newtx defines \Bbbk, \openbox, \proof and \endproof; amssymb and amsthm
% then redefine them and abort. Freeing the four names before those two
% packages load is the standard fix, and it costs nothing here: we use
% amsthm only for \newtheorem.
\let\Bbbk\relax
\usepackage{amsmath,amssymb}
\let\openbox\relax
\let\proof\relax
\let\endproof\relax
\usepackage{amsthm}
\usepackage{booktabs}
\usepackage{enumitem}
\usepackage{xcolor}
\usepackage{tcolorbox}
\usepackage{titlesec}
\usepackage{graphicx}
\usepackage[hidelinks]{hyperref}
\tcbuselibrary{skins,breakable}

% ---- the course palette, from assets/css/custom.css -------------------
\definecolor{aimlprimary}{HTML}{0B3D62}   % deep blue  - structure
\definecolor{aimlaccent} {HTML}{C0392B}   % brick red  - failures
\definecolor{aimlsuccess}{HTML}{14663A}   % green      - repairs
\definecolor{aimlmath}   {HTML}{6C3483}   % purple     - the derivation
\definecolor{aimlmuted}  {HTML}{4B5563}
\definecolor{aimlrule}   {HTML}{B0BCC7}
\definecolor{aimlsoft}   {HTML}{F4F7F9}

\usepackage[a4paper,margin=32mm,top=30mm,bottom=30mm]{geometry}
\setlength{\parskip}{0.45\baselineskip}
\setlength{\parindent}{0pt}
\linespread{1.06}

% ---- headings --------------------------------------------------------
\titleformat{\section}{\Large\bfseries\color{aimlprimary}}{\thesection}{0.7em}{}
\titleformat{\subsection}{\large\bfseries\color{aimlprimary}}{\thesubsection}{0.7em}{}
\titleformat{\subsubsection}{\bfseries\color{aimlmuted}}{}{0em}{}
\titlespacing*{\section}{0pt}{1.6\baselineskip}{0.5\baselineskip}
\titlespacing*{\subsection}{0pt}{1.2\baselineskip}{0.35\baselineskip}

% ---- boxes -----------------------------------------------------------
% One box per role, and the role is the colour. A student who has seen the
% slides already knows what each one means before reading a word of it.
\newtcolorbox{keybox}[1][]{
  breakable, enhanced, colback=aimlsoft, colframe=aimlprimary,
  boxrule=0.4pt, left=8pt, right=8pt, top=6pt, bottom=6pt,
  fonttitle=\bfseries\small, coltitle=white, colbacktitle=aimlprimary,
  attach boxed title to top left={yshift=-2mm, xshift=4mm},
  boxed title style={boxrule=0pt, arc=1pt}, #1}

\newtcolorbox{failbox}[1][]{
  breakable, enhanced, colback=white, colframe=aimlaccent,
  boxrule=0.9pt, left=8pt, right=8pt, top=6pt, bottom=6pt,
  fonttitle=\bfseries\small, coltitle=white, colbacktitle=aimlaccent,
  attach boxed title to top left={yshift=-2mm, xshift=4mm},
  boxed title style={boxrule=0pt, arc=1pt}, #1}

\newtcolorbox{derivbox}[1][]{
  breakable, enhanced, colback=white, colframe=aimlmath,
  boxrule=0.6pt, left=8pt, right=8pt, top=6pt, bottom=6pt,
  fonttitle=\bfseries\small, coltitle=white, colbacktitle=aimlmath,
  attach boxed title to top left={yshift=-2mm, xshift=4mm},
  boxed title style={boxrule=0pt, arc=1pt}, #1}

% An exercise. Deliberately the same shape as the notebook's `try` field:
% one change, and what should happen to the output.
\newtcolorbox{trybox}[1][]{
  breakable, enhanced, colback=white, colframe=aimlsuccess,
  boxrule=0.5pt, left=8pt, right=8pt, top=6pt, bottom=6pt,
  fonttitle=\bfseries\small, coltitle=white, colbacktitle=aimlsuccess,
  attach boxed title to top left={yshift=-2mm, xshift=4mm},
  boxed title style={boxrule=0pt, arc=1pt}, #1}

% ---- small semantics -------------------------------------------------
\newcommand{\scope}[1]{\textcolor{aimlmuted}{\small #1}}
\newcommand{\term}[1]{\textbf{#1}}
\newcommand{\code}[1]{\texttt{\small #1}}
\newcommand{\lecture}[1]{Lecture~#1}

\newtheorem{definition}{Definition}[section]
\newtheorem{proposition}[definition]{Proposition}

% ---- title block -----------------------------------------------------
% \notestitle{19}{Information retrieval: the lexical foundation}{SciFact (BEIR)}
\newcommand{\notestitle}[3]{%
  \thispagestyle{empty}
  \begin{flushleft}
    {\color{aimlmuted}\small\sffamily
     APPLICATIONS OF MACHINE LEARNING \quad$\cdot$\quad
     BSc Mathematics of Artificial Intelligence}\\[2mm]
    {\color{aimlrule}\rule{\textwidth}{0.8pt}}\\[4mm]
    {\color{aimlmuted}\large Lecture #1 \quad$\cdot$\quad Part V \quad$\cdot$\quad
     Extended notes}\\[3mm]
    {\Huge\bfseries\color{aimlprimary}\baselineskip=1.15\baselineskip #2\par}\vspace{4mm}
    {\color{aimlmuted}\large Dataset: #3}\\[3mm]
    {\color{aimlrule}\rule{\textwidth}{0.8pt}}
  \end{flushleft}
  \vspace{2mm}
  \begin{keybox}[title={Why these notes exist}]
    Lectures 19--22 sit outside G\'eron, the course's primary text. For those
    four lectures \emph{these notes are the primary source}, and they are
    examined on exactly the same terms as the textbook parts. Everything in
    the slide deck for this lecture is here, with the arguments written out at
    length rather than compressed onto a slide; the numbers are the ones the
    accompanying notebook prints, and every one of them is a measurement on
    the corpus named above rather than a figure quoted from a paper.
  \end{keybox}
  \vspace{1mm}
  {\small\tableofcontents}
  \vspace{2mm}
  \noindent{\color{aimlrule}\rule{\textwidth}{0.4pt}}
  \vspace{1mm}
}


\begin{document}
\notestitlefor{5}{I}{Regularisation and the bias--variance trade-off}{Titanic, 891 passengers}{Chapter 4}

\section{Two questions the sweep cannot answer}

Lecture 4 fitted a logistic regression to 712 passengers and then added
polynomial capacity, recording both curves at every degree.

\begin{center}
\small
\begin{tabular}{rrrr}
\toprule
\textbf{Degree} & \textbf{Columns} & \textbf{Training log loss} &
\textbf{Held-out log loss} \\
\midrule
1 &  22 & 0.395 & 0.468 \\
2 &  32 & 0.381 & 0.467 \\
3 &  52 & 0.355 & 0.538 \\
5 & 143 & 0.286 & 1.922 \\
\bottomrule
\end{tabular}
\end{center}

One curve falls forever; the other turns and climbs. Two questions follow, and
they are different questions. \emph{The held-out error is 0.468 and the
training error is 0.395 --- where does the difference come from?} And
\emph{0.468 against an anchor of 0.666 --- how far down could it ever go?}

The two have different repairs, and the repairs point in opposite directions.
Attack the wrong one and you spend an afternoon making the model worse.

\section{Derivation: the bias--variance decomposition}

\subsection{First, what is random}

Not the passenger. Fix one passenger $x$ and ask about them repeatedly. What
varies is the \emph{training set} $D$ --- you happened to draw these 712 rows
--- and the \emph{label} $y$, because two people with identical recorded values
did not both survive. Every expectation below is over those two sources and
nothing else; the model is deterministic given $D$.

\begin{center}
\small
\begin{tabular}{llc}
\toprule
\textbf{Symbol} & \textbf{Meaning} & \textbf{Observed?} \\
\midrule
$\hat p_D(x)$ & what your model predicts, having seen $D$ & yes \\
$\bar p(x) = \mathbb{E}_D[\hat p_D(x)]$ & the average prediction over all
                                          training sets & no --- estimated \\
$p^{*}(x)$ & the true probability that this passenger survives & never \\
\bottomrule
\end{tabular}
\end{center}

$y \in \{0,1\}$ is what actually happened. It is never $p^{*}$, and that gap
becomes the third term.

\subsection{The decomposition}

\begin{derivbox}[title={Add and subtract, twice}]
Hold $y$ fixed and average over $D$ only:
\[ \bigl(y - \hat p_D\bigr)^2 = \Bigl(
   \underbrace{y - \bar p}_{\text{does not depend on } D}
   + \underbrace{\bar p - \hat p_D}_{\text{mean } 0}\Bigr)^2. \]
The cross term is $2(y - \bar p)\,\mathbb{E}_D[\bar p - \hat p_D] = 0$, because
$\bar p$ is \emph{defined} as $\mathbb{E}_D[\hat p_D]$. So
\[ \mathbb{E}_D\bigl[(y - \hat p_D)^2\bigr]
   = (y - \bar p)^2
   + \underbrace{\mathbb{E}_D\bigl[(\hat p_D - \bar p)^2\bigr]}_{\textbf{variance}}. \]
Repeat the trick on $(y - \bar p)^2$, inserting $p^{*}$. Since
$\mathbb{E}_y[y] = p^{*}$ the cross term vanishes again, and
$\operatorname{Var}(y) = p^{*}(1-p^{*})$ because $y$ is Bernoulli:
\[ \mathbb{E}_y\bigl[(y-\bar p)^2\bigr]
   = \underbrace{(p^{*} - \bar p)^2}_{\textbf{squared bias}}
   + \underbrace{p^{*}(1-p^{*})}_{\textbf{noise}}. \]
\end{derivbox}

\[ \mathbb{E}\bigl[(y - \hat p_D)^2\bigr]
   = (p^{*} - \bar p)^2
   + \mathbb{E}_D\bigl[(\hat p_D - \bar p)^2\bigr]
   + p^{*}(1 - p^{*}). \]

\begin{center}
\small
\begin{tabular}{lp{0.40\textwidth}l}
\toprule
\textbf{Term} & \textbf{Means} & \textbf{Cured by} \\
\midrule
squared bias & wrong on average --- the model family cannot reach the truth &
               more capacity \\
variance     & unstable --- a different sample gives a different answer &
               less capacity, or more data \\
noise        & the label is not a function of the recorded columns & nothing \\
\bottomrule
\end{tabular}
\end{center}

Generally the identity is $\mathbb{E}[(y-\hat f)^2] = \text{bias}^2 +
\text{variance} + \operatorname{Var}(y \mid x)$; $p^{*}(1-p^{*})$ is that last
term for a Bernoulli label, one instance of the identity rather than the
identity itself.

The first two repairs point in opposite directions. That is the whole
difficulty, and it is why it is called a trade-off.

\begin{failbox}[title={Failure condition --- an honest caveat before we measure anything}]
The identity is exact \emph{for squared error}. Our headline metric is log
loss, and log loss does not split into three additive terms like this.

So the decomposition is measured on the \emph{Brier score} --- the squared
error of a probability, which Lecture 4 already reported alongside log loss.
The shape of the answer transfers; the numbers are Brier numbers and are
labelled as such. Saying ``bias--variance'' about a log-loss curve is common
and slightly wrong. Doing it knowingly, and saying which metric you
decomposed, is the difference.
\end{failbox}

\section{Measuring all three terms}

\subsection{The term you cannot fix}

$p^{*}$ is never observed --- but where several passengers share exactly the
same recorded values, every model of those columns \emph{must} give them the
same number. Whatever their outcomes do inside that cell is a floor.

Grouping on sex, class, age band, family band and port gives 133 distinct
cells, of which 102 hold at least two passengers, covering 858 people. Using
$k(m-k)/(m(m-1))$, the unbiased estimator of $p(1-p)$ from $m$ Bernoulli draws,
the floor is \textbf{0.121} Brier --- measured, not assumed.

56 of the 102 shared cells are mixed: 657 passengers, 74\% of everyone on
board, sit in a cell whose answer cannot be certain. The starkest is 62
third-class men aged 26 to 40 travelling alone from Southampton, with every
recorded value identical; 13 of them survived, and no model of these columns
can separate them.

It is an \emph{under}-estimate: two passengers in the same cell may still
differ in ways the manifest never recorded. The real floor is at least this
high.

\subsection{The other two, by resampling}

Draw 200 training sets of 400 rows each from the 712 we hold, fit the pipeline
on each, and predict the same 179 test passengers. That is the expectation over
$D$, done by brute force --- the only way to see a quantity defined over
training sets you did not draw. Bias and noise cannot be separated without
$p^{*}$, so they are reported as one column, and the floor above is what splits
them.

\begin{keybox}[title={Check the identity before believing the plot}]
\begin{verbatim}
>>> abs(total.mean() - variance.mean() - bias2_pn.mean())
2.7755575615628914e-17
\end{verbatim}
Machine epsilon on a double. The three columns below are not a \emph{model} of
the error --- they are the error, rearranged.

Do this every time you implement a decomposition. If the residual is not at
floating-point noise you have a bug, and the picture will still look plausible.
\end{keybox}

\begin{center}
\small
\begin{tabular}{rrrr}
\toprule
\textbf{Degree} & \textbf{bias$^2$ + noise} & \textbf{variance} &
\textbf{total} \\
\midrule
1 & 0.130 & 0.006 & 0.136 \\
2 & 0.134 & 0.011 & 0.146 \\
3 & 0.138 & 0.029 & 0.167 \\
4 & 0.163 & 0.059 & 0.222 \\
5 & 0.181 & 0.068 & 0.248 \\
6 & 0.182 & 0.071 & 0.253 \\
\bottomrule
\end{tabular}
\end{center}

Read the first column: 0.130 at degree 1, against a measured noise floor of
0.121. Almost all of it is noise, so the squared bias of a degree-1 logistic
model on these features is under 0.01.

\begin{keybox}[title={There was never much bias to buy back}]
Every extra degree bought variance and nothing else --- the variance term is
multiplied by twelve between degree 1 and degree 6 while everything else is
nearly flat. That one sentence decides the rest of the lecture.
\end{keybox}

\subsection{Putting it back on the two curves}

The training curve falls forever because it is measured on the rows that were
fitted, so it sees no variance term at all. The held-out curve is the sum of
all three terms, and from degree 3 the variance term dominates.

\begin{keybox}[title={The identification, stated once}]
The vertical gap between the two curves \emph{is} the variance term. At degree
1 the gap is 0.073; at degree 5 it is 1.636. The gap grew by a factor of 22
while nothing about the data changed.
\end{keybox}

There are only two shapes, and telling them apart takes ten seconds. Both
curves plateau high and together: \emph{bias} --- the model is too rigid, and
more data changes nothing. A gap still open at the last point: \emph{variance}
--- more data would help, and you may not be able to get any. The repairs are
opposites, and applying the left-hand repair to a right-hand curve is the most
common wasted afternoon in applied machine learning.

\subsection{Learning curves, and a knob we do not have}

A learning curve puts \emph{rows} on the $x$-axis and answers a question the
degree sweep cannot: would more data help?

\begin{center}
\small
\begin{tabular}{lrr}
\toprule
& \textbf{Degree 1} & \textbf{Degree 5} \\
\midrule
Held-out log loss at 76 rows  & 2.853 & 9.889 \\
Training log loss at 640 rows & 0.395 & 0.286 \\
Held-out log loss at 640 rows & 0.474 & 1.911 \\
Gap at 640 rows               & 0.078 & 1.625 \\
\bottomrule
\end{tabular}
\end{center}

The degree-5 held-out curve falls from 9.9 to 1.9 as rows arrive and is still
falling steeply at 640. Extrapolating, that model would need several times our
dataset before becoming competitive.

\begin{failbox}[title={There are 891 passengers on the Titanic}]
There will never be more. Rows are not a knob we can turn, so the only repair
left is to shrink the model.
\end{failbox}

A \emph{validation} curve puts one hyperparameter on the $x$-axis instead, and
answers which setting is right and whether the optimum is interior. Read the
ends first: a minimum sitting on the edge of the range means the range was too
narrow, and the answer you read off it is an artefact of where you stopped
looking.

\section{Three faults at degree 5}

Degree 5 scored 1.922 --- almost three times worse than \emph{saying nothing at
all}. Not worse than the best model; worse than the anchor of 0.666, worse than
a model that has learned only how many people died.

How is that possible when held-out accuracy was 76.0\%? Because
$-\log(0.01) \approx 4.6$: a single passenger given probability 0.01 who then
survives contributes 4.6 to the average on its own. Accuracy records that
passenger as one error and moves on. Log loss is unbounded above --- not a
defect, but the metric saying that confident-and-wrong is a different failure
from wrong.

\subsection{Fault 1 --- variance}

\begin{center}
\small
\begin{tabular}{lrr}
\toprule
& \textbf{Degree 1} & \textbf{Degree 5} \\
\midrule
Columns              & 22 & 143 \\
Training rows        & 712 & 712 \\
Rows per weight      & 32 & 5 \\
Variance term, Brier & 0.006 & 0.068 \\
\bottomrule
\end{tabular}
\end{center}

Five rows per weight. Move one passenger between folds and the fitted surface
moves with them. This is the diagnosis the derivation was built to deliver ---
and two more faults sit underneath it.

\subsection{Fault 2 --- the minimiser does not exist}

If some $\boldsymbol\theta$ separates the classes perfectly in the expanded
space, then $\sigma(c\,\boldsymbol\theta^{\intercal}\mathbf{x})$ tends to the
right label for every row as $c \to \infty$, so the training log loss falls
monotonically towards zero as $\lVert\boldsymbol\theta\rVert \to \infty$. The
infimum is never attained.

\begin{center}
\small
\begin{tabular}{lrrrrrr}
\toprule
\textbf{Degree} & 1 & 2 & 3 & 4 & 5 & 6 \\
\midrule
Converged & yes & yes & yes & no & no & no \\
Largest $|\theta_j|$ & 4.87 & 5.43 & 6.32 & 18.70 & 11.07 & 3.98 \\
\bottomrule
\end{tabular}
\end{center}

\subsection{Fault 3 --- the minimiser is not unique}

\code{FamilySize = SibSp + Parch + 1} exactly, on every row. Standardise the
five numeric columns, form $\mathbf{X}^{\intercal}\mathbf{X}$ and look at its
spectrum:

\begin{center}
\small
\begin{verbatim}
>>> np.linalg.eigvalsh(Z.T @ Z)
array([1.81e-12, 4.29e+02, 5.43e+02, 7.90e+02, 1.80e+03])
\end{verbatim}
\end{center}

The smallest eigenvalue is zero to the precision the arithmetic allows: four
columns of information in five columns of data, condition number about
$9\times10^{14}$. A double carries about $10^{16}$ of relative precision, so
roughly one significant digit survives the solve.

\section{One term repairs all three}

\subsection{What adding $\alpha\mathbf{I}$ does to the spectrum}

\begin{derivbox}[title={Three lines, and Lecture 2 is closed}]
If $\mathbf{X}^{\intercal}\mathbf{X}\mathbf{v} = \lambda\mathbf{v}$ then
\[ \bigl(\mathbf{X}^{\intercal}\mathbf{X} + \alpha\mathbf{I}\bigr)\mathbf{v}
   = (\lambda + \alpha)\mathbf{v}: \]
same eigenvectors, every eigenvalue moved up by exactly $\alpha$.
$\mathbf{X}^{\intercal}\mathbf{X}$ is positive semi-definite, so $\lambda \ge
0$, so $\lambda + \alpha \ge \alpha > 0$. No eigenvalue is zero, so the matrix
is invertible \emph{for every} $\alpha > 0$.
\end{derivbox}

Lecture 2 derived $\hat{\boldsymbol\theta} = (\mathbf{X}^{\intercal}
\mathbf{X})^{-1}\mathbf{X}^{\intercal}y$ and left open what to do when that
inverse does not exist. Those three lines are the answer.

It also bounds the conditioning:
\[ \kappa\bigl(\mathbf{X}^{\intercal}\mathbf{X} + \alpha\mathbf{I}\bigr)
   = \frac{\lambda_{\max} + \alpha}{\lambda_{\min} + \alpha}
   \le \frac{\lambda_{\max} + \alpha}{\alpha}, \]
and the bound on the right does not mention $\lambda_{\min}$ at all. Whatever
the data do --- exact collinearity included --- the conditioning is controlled
by a number you choose. At $\alpha = 1$ the condition number of our singular
matrix falls from $9\times10^{14}$ to 1798. The dependence is not repaired; it
is \emph{dominated}. And since $\kappa$ set the step count for gradient descent
in Lecture 4, the penalty is not only a statistical device --- it is what makes
the optimisation tractable.

\subsection{The idea, and the three penalties}

Stop asking for the best fit. Ask for the best fit \emph{among small models}:
\[ J(\boldsymbol\theta) = \underbrace{L(\boldsymbol\theta)}_{\text{fit the data}}
   + \underbrace{\alpha\,\Omega(\boldsymbol\theta)}_{\text{stay small}}. \]
Setting $\alpha = 0$ recovers Lecture 4 exactly.

\term{Ridge} uses $\Omega = \tfrac12\lVert\mathbf{w}\rVert_2^2$, where
$\mathbf{w}$ is $\boldsymbol\theta$ without the bias --- the intercept is
\emph{not} penalised, because shrinking it would pull every prediction towards
zero rather than towards the mean. On the log loss the objective becomes
\emph{coercive}: as $\lVert\mathbf{w}\rVert \to \infty$ the penalty tends to
infinity while the log loss is bounded below by zero, so a minimiser exists,
and strict convexity makes it unique. Faults 2 and 3 are both gone, and we have
not looked at the data once.

\term{Lasso} uses $\sum_j|\theta_j|$ instead, and the word \emph{selection} in
its name is the point.

\begin{derivbox}[title={Why $\ell_1$ produces exact zeros and $\ell_2$ does not}]
Look at the gradient of the penalty near $\theta_j = 0$. For ridge,
$\partial(\theta_j^2)/\partial\theta_j = 2\theta_j \to 0$: the push towards zero
fades as you approach it, so you arrive only in the limit. For lasso,
$\partial|\theta_j|/\partial\theta_j = \pm 1$: the push keeps its size right up
to zero, so if the data pulls with less than $\alpha$ the weight reaches zero
and stops.

At $\theta_j = 0$ the absolute value is not differentiable; the subgradient is
the whole interval $[-1,1]$, and any data pull inside that interval is held at
exactly zero.
\end{derivbox}

\term{Elastic net} mixes the two with a ratio $r$, and is preferable to lasso
when the features are correlated or outnumber the instances, because lasso
behaves erratically there. One of those holds here: $\text{Age}$ and
$\text{Age}^2$ are about as correlated as two columns get.

\begin{failbox}[title={Failure condition --- an unscaled penalty}]
The penalty $\sum\theta_j^2$ treats every weight identically. Rescale a column
by $c$ and its fitted weight scales by $1/c$, so its contribution to the
penalty scales by $1/c^2$.

Changing the \emph{unit} of a column therefore changes which coefficients ridge
decides to shrink, and nothing in the output says so. \code{StandardScaler} is
already inside the pipeline; a penalty is the reason it is not optional.
\end{failbox}

And a trap in the API: \code{LogisticRegression} takes $C = 1/\alpha$, not
$\alpha$. Small \code{C} is strong regularisation. The default is
\code{C=1.0} with \code{penalty="l2"}, so scikit-learn's logistic regression is
regularised unless you say otherwise --- Lecture 4 switched it off on purpose,
and that choice is what left three faults to repair.

\subsection{The sweep, on the model that failed}

\begin{center}
\small
\begin{tabular}{lrrr}
\toprule
\textbf{Penalty at degree 5} & \textbf{Best C} & \textbf{Log loss} &
\textbf{Non-zero weights} \\
\midrule
none --- Lecture 4 & --- & 1.922 & 143 \\
ridge & 0.01 & 0.706 & 143 \\
elastic net, $r = 0.5$ & 0.0001 & 0.684 & 3 \\
lasso & 0.00032 & 0.683 & 3 \\
\bottomrule
\end{tabular}
\end{center}

Ridge alone removes 1.215 of log loss. Lasso removes 1.239 and throws away 140
of the 143 columns while doing it. Both minima sit near the strong-penalty end
of the range, so widen the grid before believing either --- and the lasso curve
is nearly flat there, so ``the best \code{C}'' is a plateau, not a point.

\subsection{And now the uncomfortable part}

\begin{center}
\small
\begin{tabular}{lr}
\toprule
\textbf{Model} & \textbf{Cross-validated log loss} \\
\midrule
Degree 5, lasso, tuned & 0.683 \\
Degree 2, no penalty at all & 0.467 \\
\bottomrule
\end{tabular}
\end{center}

The best regularised degree-5 model does not beat the plain degree-2 model we
already had. Regularisation repaired an over-complex model; it did not turn it
into a better one than the simple model --- because there was never much bias
to buy back. The decomposition said so twenty minutes earlier, and this is what
that sentence looks like in a table.

\begin{keybox}[title={So what is a penalty actually for?}]
It makes the problem \emph{well-posed}: a unique minimiser exists for every
$\alpha > 0$, whatever the columns do. It makes capacity a \emph{continuous}
dial, where degree is an integer. It is \emph{insurance} against not knowing
the right capacity in advance, which is the usual case. And it \emph{selects}
--- lasso's three surviving columns are a hypothesis about which features
matter, generated by the fit.

What it is not is a way to make an over-parameterised model beat a well-chosen
one on a dataset this small.
\end{keybox}

\subsection{Early stopping}

Fit iteratively and stop before the optimiser gets where it is going. Descent
starts at $\boldsymbol\theta = \mathbf{0}$ --- the smallest possible model ---
and each epoch moves the weights outward, so held-out error falls, reaches a
minimum, and then rises. No penalty term and no extra term in the objective:
just a stopping rule, whose knob is the epoch count.

\begin{center}
\small
\begin{tabular}{lrr}
\toprule
& \textbf{Epoch 105} & \textbf{Epoch 500} \\
\midrule
Training log loss   & 2.458 & 2.219 \\
Validation log loss & 2.371 & 3.645 \\
\bottomrule
\end{tabular}
\end{center}

Epoch 105 is the best model this run ever holds, and running to 500 costs 1.274
of log loss without ever announcing it. Read the table honestly, though: these
are bad numbers in absolute terms. Plain SGD at a constant learning rate on 143
correlated columns is a poor optimiser, and this shows the \emph{shape}, not a
competitive model. The training curve is also noisy, so ``stop at the minimum''
needs a patience rule in practice. And the stopping epoch is chosen by looking
at held-out data, so it must be chosen on validation rows, never on the test
set.

\begin{center}
\small
\begin{tabular}{llc}
\toprule
\textbf{Repair} & \textbf{Keeps $\lVert\boldsymbol\theta\rVert$ small by} &
\textbf{Zeros?} \\
\midrule
Ridge          & penalising $\sum\theta_j^2$ & no \\
Lasso          & penalising $\sum|\theta_j|$ & yes \\
Elastic net    & both, mixed by $r$ & yes \\
Early stopping & starting at zero and not going far & no \\
\bottomrule
\end{tabular}
\end{center}

All four reduce variance by shrinking the set of models the fit may reach. All
four increase bias. Which is why all four need $\alpha$ --- or an epoch count
--- chosen by cross-validation.

\section{Choosing $\alpha$, and the test set once}

\begin{failbox}[title={Failure condition --- a hyperparameter chosen on the rows you then report}]
Looping over 17 values of \code{C}, scoring each on the test set and keeping
the best, produces a minimum over seventeen \emph{noisy} estimates --- and the
minimum of noisy estimates is biased downward even when every estimate is
individually unbiased.

The code never writes to the test set. It reads it seventeen times, and reading
is enough.
\end{failbox}

Measured over 20 splits, against the honest version that picks \code{C} by
5-fold cross-validation inside the training part:

\begin{center}
\small
\begin{tabular}{lrr}
\toprule
\textbf{Reported number} & \textbf{Log loss} & \textbf{Spread over 20 splits} \\
\midrule
Chosen on the test set, scored on the test set & 0.451 & 0.049 \\
Chosen by cross-validation, scored on the test set & 0.470 & 0.045 \\
\bottomrule
\end{tabular}
\end{center}

Optimism of 0.019 of log loss, about 4.1\%. On 13 splits of 20 the first number
is the flattering one, and the two procedures picked the same \code{C} on 7
splits of 20.

Read that result in both directions. It is \emph{small} --- 0.019 against a
split-to-split spread of 0.045, invisible on one split. And it is \emph{real
and one-sided}: it never averages away, and it grows with the number of
candidates you try. Seventeen candidates is a small search; search a thousand,
which a randomised search does in an afternoon, and the same procedure hands
you a number that is confidently wrong. The size of the error scales with how
hard you looked.

The procedure without the problem is \code{GridSearchCV}: same seventeen
candidates, same runtime, \emph{two} numbers instead of one. Report both. The
second is allowed to be worse than the first, and if it is much worse that is
itself the finding.

\subsection{The honest numbers}

179 passengers, untouched since the split in Lecture 4. Five candidates, all of
whose hyperparameters were fixed before this cell ran.

\begin{center}
\small
\begin{tabular}{lrrr}
\toprule
\textbf{Model} & \textbf{Log loss} & \textbf{Brier} & \textbf{Accuracy} \\
\midrule
Degree 5, no penalty --- where the sweep ended & 1.727 & 0.189 & 74.9\% \\
Degree 5, ridge, tuned    & 0.804 & 0.173 & 77.1\% \\
Degree 5, lasso, tuned    & 0.677 & 0.244 & 65.9\% \\
Degree 1, scikit-learn defaults & 0.433 & 0.134 & 82.7\% \\
Degree 2, no penalty --- the winner & 0.427 & 0.134 & 82.7\% \\
\bottomrule
\end{tabular}
\end{center}

The winner improves on the degree-5 model Lecture 4 finished with by 1.300 of
log loss. And row four deserves as much attention: \code{LogisticRegression()}
with every default left alone --- degree 1, \code{C=1.0}, \code{penalty="l2"}
--- scores 0.433, statistically the same model as the winner.

\begin{failbox}[title={The result nobody wanted}]
Ninety minutes of ridge, lasso, elastic net and early stopping, and the best
system on the test set is a two-line model anyone could have written before
Lecture 4 started.

Report it anyway. The alternative is choosing the interesting model over the
better one.
\end{failbox}

Then what was any of it for? You now know \emph{why} degree 2 wins, and can say
it in three terms rather than as a preference. You know the floor is 0.121
Brier and that no model of these columns goes below it. You know the degree-5
model was not merely worse but \emph{ill-posed}, and exactly which term repairs
that. And you have four repairs, measured, for the next dataset --- where 143
columns will arrive with a hundred times as many rows and the answer will be
different.

A negative result you can explain is a result. A positive result you cannot
explain is a liability.

\subsection{Is 0.427 good?}

Cross-validated, the model beats the anchor by 0.199 of log loss and the
majority class by 20.5 points of accuracy; both are real. Brier 0.134 against a
measured floor of about 0.121 leaves roughly 0.013 of Brier score on the table,
in total, for \emph{any} model of these columns. And there are 179 test
passengers, so a difference of a point or two of accuracy is three or four
people and is not a difference.

That last sentence belongs in the report. A test set of 179 cannot support the
comparisons people will want to make with it.

\begin{keybox}[title={Regularisation checklist}]
\begin{enumerate}[leftmargin=1.6em,itemsep=1pt]
  \item Is the data scaled, inside the pipeline, before a penalty is applied?
  \item Is the bias term excluded from the penalty?
  \item Was $\alpha$ (or \code{C}) chosen on validation folds, never on the
        test set?
  \item If the search sat at the edge of the grid, was the grid widened?
  \item Is the \emph{unregularised} model in the comparison table, so the
        penalty has to earn its place?
  \item Was the stopping epoch, if any, chosen on held-out rows?
\end{enumerate}
The fifth is the one people skip, and it is the one that would have saved an
hour today.
\end{keybox}

\section{The notebook}

\begin{trybox}[title={What to change}]
In the penalty sweep, widen \code{np.logspace(-4, 4, 17)} to
\code{(-8, 4, 25)}. Does the lasso minimum move? If it does, the number in
\S6.3 was an artefact of where we stopped looking --- and you have just done
checklist item four on the course's own result.
\end{trybox}

\section{Exercises}

\begin{enumerate}[leftmargin=1.6em,itemsep=4pt]
  \item Write the bias--variance decomposition of the expected squared error,
        and name the one term no model can reduce. \hfill \scope{[4]}
  \item A learning curve shows training and validation error both high and
        close together, and flat as data is added. Diagnose it, and say whether
        more data would help. \hfill \scope{[5]}
  \item Ridge adds $\alpha I$ to $X^{\mathsf T}X$ before inverting. Give the
        numerical consequence, in terms of the condition number.
        \hfill \scope{[4]}
  \item Why must features be scaled before ridge or lasso, when they need not
        be for ordinary least squares? \hfill \scope{[4]}
  \item You tune $\alpha$ over a grid using cross-validation, then report the
        best cross-validated score as your estimate of future performance. Name
        the error, and quantify its direction. \hfill \scope{[5]}
\end{enumerate}

\section*{Summary}
\addcontentsline{toc}{section}{Summary}

Expected squared error splits exactly into squared bias, variance and
irreducible noise --- over the training set and the label, never over the test
point. The gap between a training curve and a held-out curve \emph{is} the
variance term; curves plateauing together mean bias instead. On this data the
noise floor is measurable at 0.121 Brier, from passengers whose recorded values
are identical, and the squared bias of a degree-1 model is under 0.01 --- so
every extra degree bought variance and nothing else.

Degree 5 failed three ways at once: high variance at five rows per weight, no
minimiser under separation, and no unique minimiser under an exact collinearity
whose smallest eigenvalue was $1.81\times10^{-12}$. Adding $\alpha\mathbf{A}$
repairs the last two for every $\alpha > 0$, which is Lecture 2's open question
answered in three lines, and bounds the condition number by a quantity you
choose.

Ridge shrinks, lasso selects, elastic net does both, early stopping does it by
not arriving. None of them beat choosing the right capacity: the winner on the
test set was degree 2 with no penalty at 0.427, against 0.683 for the best
tuned degree-5 model. A penalty is a repair, not an upgrade. And a
hyperparameter chosen on the test set cost 0.019 of log loss here, one-sided,
growing with the size of the search.

\end{document}
